\section{Analysis}
\label{sec:theoretical}

% We analyze the relationship between the loss weights $\alpha$ and $\beta$, corresponding to the competing tasks of predicting $y$ and reconstructing $x$, respectively, in our generalized unsupervised invariance induction framework. We then discuss the equilibrium of the minimax game in our adversarial instantiation. Finally, we use the results of these two analyses to provide a systematic way for tuning the loss weights $\alpha$ $\beta$, and $\gamma$. The following analyses are done assuming a model with infinite capacity.

% The following analyses are done assuming a model with infinite capacity.

\textbf{Competition between prediction and reconstruction. \quad} The prediction and reconstruction tasks in our framework are designed to compete with each other. Thus, $\eta = \frac{\alpha}{\beta}$ influences which task has higher priority in the overall objective. We analyze the affect of $\eta$ on the behavior of our framework at optimality, considering perfect disentanglement of $e_1$ and $e_2$. There are two asymptotic scenarios with respect to $\eta$: (1) $\eta \rightarrow \infty$ and (2) $\eta \rightarrow 0$. In case (1), our framework reduces to a predictor model, where the reconstruction task is completely disregarded. Only the branch $x \dashedrightarrow e_1 \dashedrightarrow y$ remains functional. Consequently, $e_1$ contains all $f \in F'$ at optimality, where $F_y \subseteq F' \subseteq F$. In contrast, case (2) reduces the framework to an autoencoder, where the prediction task is completely disregarded, and only the branch $x \dashedrightarrow e_2 \dashedrightarrow x'$ remains functional because the other input to $Dec$, $\psi(e_1)$, is noisy. Thus, $e_2$ contains all $f \in F$ and $e_1$ contains nothing at optimality, under perfect disentanglement. In transition from case (1) to case (2), by keeping $\alpha$ fixed and increasing $\beta$, the reconstruction loss starts contributing more to the overall objective, thus inducing more competition between the two tasks. As $\eta$ is gradually decreased, $f \in (F' \smallsetminus F_y) \subseteq \overline{F}_y$ migrate from $e_1$ to $e_2$ because $f \in \overline{F}_y$ are irrelevant to the prediction task but can improve reconstruction by being more readily available to $Dec$ through $e_2$ instead of $\psi(e_1)$. After a point, further decreasing $\eta$ is, however, detrimental to the prediction task as the reconstruction task starts dominating the overall objective and pulling $f \in F_y$ from $e_1$ to $e_2$.

\textbf{Equilibrium analysis of adversarial instantiation.\quad} The disentanglement and prediction objectives in our adversarial model design can simultaneously reach an optimum where $e_1$ contains $F_y$ and $e_2$ contains $\overline{F}_y$. Hence, the minimax objective in our method has a \emph{win-win equilibrium}.

\textbf{Selecting loss weights. \quad} Using the above analyses, any $\gamma$ that successfully disentangles $e_1$ and $e_2$ should be sufficient. On the other hand, $\alpha$ and $\beta$ can be selected by starting with $\alpha \gg \beta$ and gradually increasing $\beta$ as long as the performance of the prediction task improves. We found $\alpha = 100$, $\beta = 0.1$ and $\gamma = 1$ to work well for all datasets on which we evaluated the proposed model.